\documentclass{article}

\usepackage{tabularx}
\usepackage{booktabs}

\title{Problem Statement and Goals\\\progname}

\author{\authname}

\date{}

\input{../Comments.text}
\input{../Common.text}

\begin{document}

\maketitle

\begin{table}[hp]
\caption{Revision History} \label{TblRevisionHistory}
\begin{tabularx}{\textwidth}{llX}
\toprule
\textbf{Date} & \textbf{Developer(s)} & \textbf{Change}\\
\midrule
Sept 22, 26 & Qusay Qadir & problem statement, environments, goals\\
Date2 & Name(s) & Description of changes\\
... & ... & ...\\
\bottomrule
\end{tabularx}
\end{table}

\section{Problem Statement}

\wss{You should check your problem statement with the
\href{https://smiths.github.io/capTemplate/Checklists/ProbState-Checklist.pdf}
{problem statement checklist}}. 

% \wss{You can change the section headings, as long as you include the required
% information.}

\subsection{Problem}

% \wss{What problem are you trying to solve?  Why is it important? Do not talk
% about how you will solve the problem, only what the problem is.  AI is rarely
% part of the problem.  The problem is ``what'' you want to do, not ``how'' you
% want to do it.  The problem statement should be written in a way that is
% understandable to a non-technical audience.}

As code and software applications become cheaper to generate by AI assistants, 
the engineering intent behind the system design, requirements, and quality expectations 
that shape the application can become lost across prompt-and-answer conversations, 
leaving development teams unable to clearly explain why the software is structured and built the way it is. 
This creates concerns for the maintainability and scalability of the application, as well as the ability to 
connect engineering decisions to the code and its test cases. There is a lack of traceability between engineering 
intent, rationale, and design decisions, making it difficult for software engineers to remain in control of the 
system's quality, acceptance criteria, and overall goals.

\subsection{Inputs and Outputs}

\wss{Characterize the problem in terms of ``high level'' inputs and outputs.  
Use abstraction so that you can avoid details.}

\subsection{Stakeholders}

\subsection{Environment}

% \wss{Hardware and Software Environment for the users.  The developer environment
% is summarized as part of the developer plan.}

Environments is the engineering integrated development environment (IDE), where they choose to edit and write the code for their software applications.
In the IDE the user must be interacting with an AI assistant, and the hardware environment should be the standard computer (desktop or laptop) that has an
IDE installed, with AI assistant plugins.

\section{Goals}

\section{Stretch Goals}

% \wss{Goals you hope to get to, but not necessary.  They are also helpful if the
% scope of the original project needs to be expanded.}


Having a UML diagram sync that is part of the intent model, with round-trip engineering where, based on criteria changes, 
the UML diagrams evolve accordingly to provide the most accurate class relationships, as well as the structure and function of each class.
We could go more granular and include method and their types, and function parmeters and their types. 
\\\\
Rationale reconstruction is a feature that allows the user to go back in time to a specific commit or pull request and understand what the requirements 
or constraints were at that point in the development cycle. This can paint a picture of how requirements, goals, quality priorities, and design constraints have evolved, and it links 
directly to the code and test cases.
\\\\
GitHub Actions integration: when new code is about to be merged into the main branch, an agent checks it against the intent model and comments on 
which parts of the model the change touches and whether they are still valid. This would ease the PR approval phase, since reviewers would not need to know what each line of code does.

\section{Custom Documents (CDs)}

\wss{For CAS 741: State whether the project is a research project. This
designation, with the approval (or request) of the instructor, can be modified
over the course of the term.}

\wss{For SE Capstone: List your custom documents (CDs).  Potential CDs include
usability testing, code walkthroughs, user documentation, formal proof,
GenderMag personas, Design Thinking, etc.  (The full list is on the course
outline and in Lecture 02.) Normally the number of CDs will be two (one per
term).  Approval of the CDs will be part of the discussion with the instructor
for approving the project. The CDs, with the approval (or request) of the
instructor, can be modified over the course of the term.}

\newpage{}

\section*{Appendix --- Reflection}

\wss{Not required for CAS 741}

\input{../Reflection.text}

\begin{enumerate}
    \item What went well while writing this deliverable? 
    \item What pain points did you experience during this deliverable, and how
    did you resolve them?
    \item How did you and your team adjust the scope of your goals to ensure
    they are suitable for a Capstone project (not overly ambitious but also of
    appropriate complexity for a senior design project)?
\end{enumerate}  

\end{document}